% [原文 p.249]
\begin{proof}[定理 6.32 的证明]
证明沿与定理 6.27 相同的思路进行。

由引理 6.2，定理 6.32 的 (1) 与 (2) 等价。由推论 6.7，只需对指数 1 的解析终端奇点的芽，即孤立 cDV 奇点的芽证明 (1) 即可。

设 $\pi:(P\in X)\to(0\in V)$ 是到非奇异解析曲线的态射，$X_0$ 是 Du Val 的，且一般纤维 $X_\xi$ 非奇异。进一步由命题 6.8 可知：若对某非奇异曲线 $W$ 的基变换 $X\times_V W\to W$ 定理 6.32 成立，则对原来的 $V$ 定理 6.32 也成立。因此可假定 $\pi$ 具有同时极小奇点消解 $h:Z\to X$。与定理 6.27 不同，这里 $h$ 在余维 1 处是同构。置 $G=h_*^{-1}D$，运行 $(K_Z+\varepsilon G)$-MMP。其结果是得到 $h':Z'\to X$，且 $G$ 在 $Z'$ 上的双有理变换 $G'$ 是 $h'$-nef 的。此时 $G'$ 是 $h'$-自由的，且对所有 $\nu\ge0$ 有 $h'_*\cO_{Z'}(\nu G')=\cO_X(\nu D)$。由此导出定理 6.32 的方法与定理 6.27 的证明相同。
\end{proof}

% [Tutor Note: 定理 6.23 与定理 6.25 合起来给出三维 MMP 的标准出发点：先作 crepant 终端化，再作 Q-分解，所得模型同时终端且 Q-分解。证明中反复使用同一套“符号幂—MMP—相对无基点性”的论证，值得与定理 6.27 的证明对读。]

\begin{remark}{6.34}
若定理 6.32 中的 $X$ 是代数簇，则与注意 6.29 一样，可以只用射影簇的 MMP 给出证明。大筋相同，故沿用注意 6.29 的记号，这里只说明需要额外补充什么。

所需要的是 $\bar Y_n\times_{\C^n}W$ 上的整体 Weil 除子 $D'$，使其在 $Y_n\times_{\C^n}W^\circ$ 上的限制与 $D$ 在局部解析同构 $X\times_{V^\circ}W^\circ\to\bar Y_n\times_{\C^n}W^\circ$ 下的像线性等价。注意 $\bar Y_n\times_{\C^n}W^\circ$ 至多具有 cDV 奇点，故由定理 5.34 与定理 5.22，它至多具有有理奇点。以 $F$ 表示 $Y_n\to\bar Y_n$ 中所关注的奇异纤维，则由 $Y_n$ 的构造（定理 4.39）知 $\operatorname{Pic}Y_n\to\operatorname{Pic}F$ 是满射（定理 4.39(10)）。于是所需的 $D'$ 由下面的引理 6.35 给出，只用射影簇的 MMP 即可完成。
\end{remark}

% [原文 p.250]
\begin{lemmap}{6.35}
设 $h:Z\to X$ 是 $\C$ 上的至多具有有理奇点的簇的奇点消解，$x\in X$ 是使限制映射 $\operatorname{Pic}Z\to\operatorname{Pic}h^{-1}(x)$ 为满射的点。设 $x\in U\subset X^{an}$ 是满足 $H^1(h^{-1}(U),\cO_{h^{-1}(U)})=0$ 的开邻域，$D$ 是 $U$ 上的 Weil 除子。此时存在 $X$ 上的代数 Weil 除子 $D'$，使 $D$ 与 $D'|_U$ 线性等价。
\end{lemmap}

\begin{proof}
注意 $X^{an}$ 至多具有有理奇点（定理 2.48）。由引理 4.13，$h_*^{-1}D|_{h^{-1}(U)}$ 由其像在 $\operatorname{Pic}h^{-1}(x)$ 中决定。因此由假设，存在 $Z$ 上的 Cartier 除子 $F$，使 $F|_{h^{-1}(U)}\sim h_*^{-1}D|_{h^{-1}(U)}$。于是 $h_*F|_U\sim D|_U$。
\end{proof}

\begin{lemmap}{6.36}
设 $X$ 只具有终端奇点，$Z\subset X$ 是余维 $\ge2$ 的闭子集。此时集合
\[
\Sigma(X,Z)=\{X\ \text{上的除子}\ E \mid \varnothing\ne\operatorname{center}_X E\subset Z,\ a(E,X,\Delta)\le1\}
\]
是有限集。
\end{lemmap}

% [Tutor Note: 引理 6.36 的证明只用 $a(E,X)$，页图与英文底本在集合定义处却写 $a(E,X,\Delta)$；引理的假设中并无 $\Delta$，此处疑为原书笔误，译文照录页图。]

% [Tutor Note: $\operatorname{center}_X E$ 指除子 $E$ 在 $X$ 上的中心，即其在 $X$ 中的像的闭包；有终端奇点时，中心落在低维子集上的除子只能有限多个，这正是定理 6.25 能对 $|\Sigma(X,\operatorname{Sing}(X))|$ 作归纳的原因。]

\begin{proof}
设 $f:Y\to X$ 是使 $\operatorname{Ex}(f)\cup f^{-1}(Z)=\bigcup_{i=1}^m E_i$ 成为 snc 除子的奇点消解。写 $f^*K_X=K_Y-\sum_i a(E_i,X)E_i$。我们来证明 $\Sigma(X,Z)\subset\{E_1,\dots,E_m\}$。为此假定 $E\in\Sigma(X,Z)$ 在 $Y$ 上是例外的。由于 $Y$ 非奇异，有 $a(E,Y)\ge1$。于是下式成立：
\[
\begin{aligned}
a(E,X)&=a\Bigl(E,Y,-\sum_i a(E_i,X)E_i\Bigr)\\
&=a(E,Y)+\sum_i a(E_i,X)\cdot\operatorname{mult}_E E_i>1.
\end{aligned}
\]
这与 $E\in\Sigma(X,Z)$ 矛盾，故命题得证。
\end{proof}

\begin{proof}[定理 6.25 的证明]
注意 $\operatorname{codim}_X\operatorname{Sing}(X)\ge3$。对 $|\Sigma(X,\operatorname{Sing}(X))|$ 归纳（引理 6.36）。若 $X$ 是 Q-分解的，则无须证明。因此不妨假定存在不是 Q-Cartier 的 Weil 除子 $D$。此时由引理 6.2 与定理 6.32，$f:Y=\Proj_X R(X,D)\to X$ 没有例外除子，且 $f_*^{-1}D$ 是 Q-Cartier 的。因此 $K_Y=f^*K_X$，且 $Y$ 是终端的。于是 $\Sigma(X,\operatorname{Sing}(X))\supset\Sigma(Y,\operatorname{Sing}(Y))$。另一方面，$\dim\operatorname{Ex}(f)=1$，爆破 $\operatorname{Ex}(f)$ 中的曲线，便得到满足 $a(E,X)=1$ 且 $\operatorname{center}_X E\subset\operatorname{Sing}(X)$ 的除子 $E$。特别地，$E\in\Sigma(X,\operatorname{Sing}(X))\setminus\Sigma(Y,\operatorname{Sing}(Y))$，故 $\Sigma(X,\operatorname{Sing}(X))\ne\Sigma(Y,\operatorname{Sing}(Y))$。由归纳法，证明完成。

% [原文 p.251]
若 $X$ 是非固有的三维解析空间且 $\operatorname{Sing}(X)=\infty$，则由上述构造，作为 $X$ 的符号爆破 $\sigma_i$ 的极限，得到固有双有理态射 $f:Y\to X$。这是因为在 $X$ 的任意紧致子集 $K$ 上，当 $i\gg0$ 时符号爆破都是同构。
\end{proof}

\section{标准翻模}
\takeaway{本节证明三维标准翻模的存在（定理 6.44、推论 6.45）与终止的一个特殊情形（定理 6.43），一般情形的终止引用 [Kol89] 的结论而不加证明。}

本节的主要结果是三维标准翻模的存在（定理 6.44、推论 6.45）与三维标准翻模终止的一个特殊情形（定理 6.43）。三维标准翻模终止的一般情形是已知的（[Kol89]），但说明会变得复杂，故本书把讨论限制在所需要的特殊情形。

\begin{remark}{6.37}
本节的讨论（只要使用相应设定下的锥定理与 MMP 等）在射影的、非固有代数的以及解析的设定下都可行。在解析的情形，书中未必明确写出，但与此前一样，把讨论限制在紧致集合的适当邻域上进行。
\end{remark}

% [Tutor Note: 三种设定可以平行处理，是因为所用的锥定理、MMP 与相对无基点性在各设定下都有对应版本；解析情形要靠紧致子集上的有限性把整体问题化为局部问题，这是本节反复提到注意 6.37 的原因。]

\begin{theoremp}[{[Kaw88, Kol89]}]{6.38}
设给定射影态射 $f_i:X_i\to Z$（$i=1,2$），以及与之相容的 $Z$ 上的双有理映射 $\pi:X_1\dashrightarrow X_2$。对 $i=1,2$，设 $X_i$ 是至多具有终端奇点的 Q-分解三维簇，且 $K_{X_i}$ 是 $f_i$-nef 的。设 $D_2$ 是 $X_2$ 上的 $f_2$-丰富有效除子，置 $D_1:=\pi_*^{-1}(D_2)$。此时 $\pi$ 等于有限个 $D_1$-翻模的合成。（解析情形参见注意 6.37。）
\end{theoremp}

\begin{proof}
从 $X_1$ 出发运行 $(K_{X_1}+\varepsilon D_1)$-MMP。由推论 3.54，$D_1$ 在 $X_1$ 的余维 2 子集之外是丰富的。因此 MMP 过程中只出现翻模。作为 MMP 的结果，得到 $X_3$、nef 的 $D_3$，以及余维 1 处同构的双有理映射 $X_3\dashrightarrow X_2$。定理 6.38 可由引理 6.39 得到。
\end{proof}

\begin{lemmap}{6.39}
设 $\pi:X_1\dashrightarrow X_2$ 是 $Z$ 上的双有理映射，且 $\pi$ 在余维 1 处是同构。
% [原文 p.252]
其次假定存在如下 $D_1,D_2$。
\begin{enumerate}[label=（\arabic*）,leftmargin=3.2em]
\item $X_2$ 上的 $f_2$-丰富 Cartier 除子 $D_2$，使 $D_1:=\pi_*^{-1}D_2$ 成为 $f_1$-nef 的 Cartier 除子。
\item $X_1$ 上的 $f_1$-丰富 Cartier 除子 $E_1$，使 $E_2:=\pi_*E_1$ 是 Cartier 的。
\end{enumerate}
此时 $\pi$ 是同构。
\end{lemmap}

\begin{proof}
对 $n\gg0$，$L_1=E_1+nD_1$ 在 $X_1$ 上是 $f_1$-丰富的（定理 1.18），$L_2=E_2+nD_2$ 在 $X_2$ 上是 $f_2$-丰富的。由 $H^0(X,\cO(mL_1))=H^0(X_2,\cO(mL_2))$，$|mL_1|$ 与 $|mL_2|$ 分别把 $X_1$ 与 $X_2$ 嵌入 $Z$ 上的同一像。因此 $\pi$ 是同构。
\end{proof}

把相对锥定理 3.25 改写成特殊情形。

\begin{propositionp}{6.40}
设 $f:X\to Z$ 是射影态射，$D$ 是 $X$ 上的有效除子且 $-D$ 是 $f$-丰富的。设 $\Delta$ 是 $X$ 上的有效 Q-除子，且 $(X,\Delta)$ 是 klt 的、$K_X+\Delta\equiv_f0$。此时有效 1-链的锥 $NE(X/Z)$（$\subset N(X/Z)\cong\mathbb{R}^{\rho(X/Z)}$）由有限个极射线生成。（解析情形参见注意 6.37。）
\end{propositionp}

\begin{proof}
对充分小的 $\varepsilon>0$，$(X,\Delta+\varepsilon D)$ 是 klt 的（推论 2.35(2)）。把相对锥定理 3.25(2) 应用于 $K_X+\Delta+\varepsilon D+(\varepsilon/2)(-D)\equiv_f(\varepsilon/2)D$，可知 $\overline{NE}(X/Z)$ 由有限个极射线与 $\overline{NE}(X/Z)_{D\ge0}$ 生成。然而 $-D$ 是 $f$-丰富的，故 $\overline{NE}(X/Z)_{D\ge0}=0$（定理 1.44）。于是得到命题 6.40。
\end{proof}

设给定两个双有理态射 $f_i:Y_i\to Z$（$i=1,2$），则 $Z$ 上的双有理映射
\[
\sigma=f_2^{-1}\circ f_1:Y_1\dashrightarrow Y_2
\]
唯一确定。当这个 $\sigma$ 是同构时，将 $f_1$ 与 $f_2$ 视为同一。

\begin{corollaryp}{6.41}
设 $Z$ 是至多具有标准奇点的射影簇或 Stein 簇，$f:X\to Z$ 是 crepant 射影双有理态射。此时 $NE(X/Z)$ 由有限个极射线生成。此外，在射影双有理态射 $h:Y\to Z$ 中，$Y$ 正规且被 $X$ 支配者至多只有有限个。
\end{corollaryp}

% [原文 p.253]
\begin{proof}
由引理 6.28，存在 $X$ 上的有效除子 $D$，使 $-D$ 是 $f$-丰富的。此时 $NE(X/Z)$ 有限生成（命题 6.40）。推论 6.41 中所述那样的任意 $Y$，与 $NE(X/Z)$ 的任意面 $NE(X/Y)$ 一一对应。
\end{proof}

\begin{theoremp}[{[KM87]}]{6.42}
设 $Z$ 是至多具有标准奇点的三维簇。此时除 $Z$-同构外，从正规簇 $X$ 出发的射影 crepant 双有理态射 $h:X\to Z$ 至多只有有限个。（解析情形参见注意 6.37。）
\end{theoremp}

\begin{proof}[证明（据 [Kol89]）]
为证有限性，只需在 $Z$ 的拟射影或 Stein 开集上讨论。因此可假定 $Z$ 是射影的或 Stein 的。定理 6.42 中所述的任意 $X$ 都具有 Q-分解的终端化，故只需证明终端且 Q-分解的 $X$ 的有限性（推论 6.41）。选取并固定那样的 $f^1:X^1\to Z$、有效丰富除子 $D^1$ 以及有效除子 $E^1\sim_{f^1}-D^1$（引理 6.28）。

于是，从任意其他 $X$ 出发，都可用有限个射线型 $D^1$-翻模的合成到达 $X^1$。由于 $D^1$-翻模的逆映射是 $E^1$-翻模，其他所有 $X$ 都可由 $X^1$ 出发施以有限个 $E^1$-翻模得到。由推论 6.41，每一步中进行 $E^1$-翻模的可能性只有有限多种。现在假定存在无限多个 $X$。此时适当选取 $E^1$-翻模 $\varphi^1:X^1\dashrightarrow X^2$，置 $E^2=\varphi^1_*E^1$，则无限多个 $X$ 可由 $X^2$ 经 $E^2$-翻模的合成得到。同样选取 $E^2$-翻模 $\varphi^2:X^2\dashrightarrow X^3$，并重复此操作。这样便得到射线型终端 $E^1$-翻模的无限列。这与射线型翻模的终止（推论 6.19）矛盾。
\end{proof}

\begin{theorem*}[三维标准翻模的终止定理（简版）]{6.43}
设 $Z$ 是至多具有标准奇点的三维簇，$f:X\to Z$ 是从正规簇出发的射影 crepant 双有理态射。设 $D$ 是有效且 Q-Cartier 的 Q-除子。此时 $Z$ 上 $D$-翻模的任意列都是有限的。（解析情形参见注意 6.37。）
\end{theorem*}

\begin{proof}
设 $X^1\dashrightarrow X^2\dashrightarrow\cdots$ 是使所有 $X^i$ 都支配 $Z$ 的 $D$-翻模的无限列。由定理 6.42，存在 $a<b$ 使 $X^a=X^b$ 且 $D^a=D^b$。然而由引理 3.38，差异系数 $a(E,X^i,\varepsilon D^i)$ 对每个 $E$ 都是非减函数，且对适当的 $E$ 有
\[
a(E,X^a,\varepsilon D^a)<a(E,X^{a+1},\varepsilon D^{a+1})\le a(E,X^b,\varepsilon D^b).
\]
这是不可能的。
\end{proof}

% [Tutor Note: 原书在定理 6.43 的陈述末印空心方块 □（通常表示证明从略或见后文），但紧接着就给出证明，与全书尾符体例不同；译文照页图保留 □，故用星号版本。]

% [原文 p.254]
\begin{theoremp}[三维标准翻模的存在 {[Kaw88, Kol89]}]{6.44}
设 $(X,0)$ 是三维标准对。此时成立：
\begin{enumerate}[label=（\arabic*）,leftmargin=3.2em]
\item 存在从 Q-分解三维簇出发的射影双有理态射 $f:\bar X\to X$，它在余维 1 处是同构。
\item 若 $e(X)>0$（定义 6.30）且 $X$ 是 Q-分解的，则存在具有下述性质的 crepant 双有理态射 $g:X'\to X$：$X'$ 是 Q-分解的，$\operatorname{Ex}(g)$ 是不可约除子，$\rho(X'/X)=1$，且 $e(X')=e(X)-1$。
\item 若 $X$ 是 Q-分解的，则对任意有效除子 $D$，存在射线型 $D$-翻模。（解析情形参见注意 6.37。）
\end{enumerate}
\end{theoremp}

\begin{proof}
对 $e(X)$ 归纳证明 (1)、(2)、(3)。这里用 $(j)_k$ 表示“当 $e(X)\le k$ 时 (j) 成立”。

若 $e(X)=0$，则 $X$ 是终端的，$(1)_0$ 无非就是 Q-分解（定理 6.25）。$(2)_0$ 是空的主张，$(3)_0$ 是终端翻模的存在（定理 6.14）。

设 $e(X)=k>0$，并假定 $(1)_{k-1}$、$(2)_{k-1}$ 与 $(3)_{k-1}$ 已证明。设 $\pi:Z\to X$ 是 $X$ 的 Q-分解终端化（定理 6.23、定理 6.25）。

注意 $\pi$ 是 crepant 的，且 $Z$ 是终端且 Q-分解的。以 $E$ 表示 $\pi$ 的所有例外除子之和。运行 $X$ 上的 $(K_Z+\varepsilon E)$-MMP。只要 $\pi^i:Z^i\to X$ 有例外除子，$K_{Z^i}+\varepsilon E^i\equiv\varepsilon E^i$ 就不是 nef 的（引理 3.39）。因此由 $e(Z^i)<k$，可以应用 $(3)_{k-1}$ 与简版终止定理 6.43。于是 MMP 在 $\pi^n:Z^n\to X$ 处终止，$Z^n$ 是 Q-分解的，且 $\pi^n$ 在余维 1 处是同构。这就证明了 $(1)_k$。

进一步假定 $X$ 是 Q-分解的。射到 Q-分解簇的态射的例外集是纯余维 1 的（推论 2.63）。因此上述 MMP 中 $\pi^n$ 是同构，而 $\pi^{n-1}:Z^{n-1}\to X$ 正是 $(2)_k$ 所要求的除子收缩。

% [原文 p.255]
设 $e(X)=k$，$h:X\to Y$ 是使 $K_X$ 为 $h$-平凡的、$-D$ 为 $h$-丰富的射线型翻模收缩。由 $(2)_k$ 得到 crepant 态射 $g:X'\to X$。置 $D':=g^*D$，从 $X'$ 出发运行 $Y$ 上的 $(K_{X'}+\varepsilon D')$-MMP。

对 $m\gg0$，$-\operatorname{Ex}(g)-mg^*D$ 是 $h\circ g$-丰富的（命题 1.45），故 $X'$ 恰有一条与 $g$ 不对应的极射线（命题 6.40）。以 $R$ 表示它。由于 $-D$ 是 $h$-丰富的，有 $(R\cdot D')<0$。从这条极射线 $R$ 出发运行 MMP，便得到如下交换图。
\[
\begin{tikzcd}
X'=X^{i_1} \arrow[r,dashed] \arrow[d,"g"] & \cdots \arrow[r,dashed] & X^{i_{i-1}} \arrow[d,"\phi^{i-1}"] \\
X \arrow[dr,"h"'] & & X^+=X^{i_i} \arrow[dl,"h^+"] \\
& Y &
\end{tikzcd}
\]
由 $e(X')=k-1$，只要 $D'$-翻模出现，就可应用 $(3)_{k-1}$。由定理 6.43，某处必出现除子收缩 $\phi^{i-1}:X^{i_{i-1}}\to X^+=X^{i_i}$。此时 $h^+:X^+\to Y$ 在余维 1 处是同构。由于 MMP 过程中出现了极射线 $R$，对适当的 $E$ 有 $a(E,X,\varepsilon D)=a(E,X',\varepsilon D')<a(E,X^+,\varepsilon D^+)$（引理 3.38）。因此 $Y$ 上 $X^+\ne X$。注意 $\operatorname{Pic}(X/Y)\otimes\Q\cong\operatorname{Pic}(X^+/Y)\otimes\Q\cong\Q$ 由 $D$、从而由 $D^+$ 生成。若 $-D^+$ 是 $h^+$-nef 的，则 $X\cong X^+$（引理 6.39），矛盾。于是 $D^+$ 是 $h^+$-丰富的，且 $X^+$ 是 $X$ 的翻模。这就证明了 $(3)_k$。
\end{proof}

\begin{corollaryp}{6.45}
设 $X$ 是至多具有标准奇点的三维簇，$f:X\to Y$ 是（不必射线型的）$D$-翻模收缩。设 $g:\bar X\to X$ 是 Q-分解，置 $\bar D=g^*D$。此时存在 $f$ 的 $D$-翻模 $f':X'\to Y$。此外，适当选取 Q-分解 $g':\bar X'\to Y$，可使诱导的双有理映射 $\bar X\dashrightarrow\bar X'$ 成为有限个射线型 $\bar D$-翻模的合成。
\end{corollaryp}

% [原文 p.256]
\begin{proof}
运行 $Y$ 上的 $(K_{\bar X}+\varepsilon\bar D)$-MMP。由于 $\bar X\to Y$ 没有例外除子，MMP 过程中只出现 $\bar D$-翻模。因此得到 $\bar X'\to Y$，使 $\bar D'$ 在 $Y$ 上是 nef 的。由相对无基点性定理 3.24，对适当的 $n>0$，$\varphi_{|n\bar D'|/Y}$ 是 $Y$ 上的态射。设 $\bar X'\to X'$ 是 $\varphi_{|n\bar D'|/Y}$ 的 Stein 分解。此时 $X'$ 就是 $f$ 的 $D$-翻模。
\end{proof}

\begin{corollaryp}[三维终端翻模的终止]{6.46}
三维终端翻模的任意列都是有限的。
\end{corollaryp}

\begin{proof}
若存在无限列，则由推论 6.45 得到三维射线型终端翻模的无限列。由推论 6.19，这不可能。
\end{proof}

% [Tutor Note: 推论 6.46 与推论 6.19 同名（原书如此）：6.19 只处理射线型终端翻模，6.46 则由之化归到一般的三维终端翻模，两者内容不同，并非重复条目。]

% [Tutor Note: 推论 6.46 把“射线型”的终止推广到一般终端翻模：先用 Q-分解把任意翻模分解为有限个射线型翻模，再引用射线型的终止性。这一“化归为射线型”的技巧在定理 6.47 的证明中同样关键。]

本书虽然不需要，但标准翻模的一般终止定理是已知的。

\begin{theorem*}[三维标准翻模的终止 {[Kol89]}]{6.47}
三维标准翻模的任意列都是有限的。（解析情形参见注意 6.37。）
\end{theorem*}

定理 6.47 按如下方法证明。首先，三维射线型标准翻模的终止，论证虽然更复杂，但可沿与定理 6.44 相同的思路证明。这样便可与推论 6.46 同样地证明即使不限于“射线型”，终止也成立。

\chapterreview{第 6 章是全书的转折点：第 3、4 章建立的锥定理、极射线收缩与差异系数理论，在这里第一次被系统用于三维双有理几何的核心操作——翻模。本章要回答三个问题：三维终端翻模与标准翻模是否存在，翻模的序列能否终止，以及终端化与 Q-分解这两个整理奇点的手术能否在 crepant 的条件下完成。这些问题此前只是以极小模型纲领应当可行的形式被提出，本章对三维情形给出了基本完整的回答（一般标准翻模的终止引用 [Kol89]）。与第 3、4 章相比，本章的结论更接近几何的最终目标：它把三维的双有理分类从纲领变成了定理。

知识脉络上，§6.1 先给出翻转与翻模的一般定义，并用引理 6.2 把翻转存在化归为分次代数 $\bigoplus\cO_Y(mB)$ 的有限生成性；§6.2 证明三维终端翻模总存在（定理 6.14），并证明三维标准翻转的终止（定理 6.17），所用不变量是 $d(X,\Delta)$；§6.3 转向终端化（定理 6.23）与 Q-分解（定理 6.25），关键技术是符号幂 $I^{(\nu)}$ 的有限生成性（定理 6.27）与符号爆破（定理 6.32）；§6.4 用 $e(X)$ 的归纳证明三维标准翻模的存在（定理 6.44、推论 6.45），并证明终端翻模的终止（推论 6.46）与标准翻模终止的简版（定理 6.43）。

核心结论可归为三条。第一，三维终端翻模保持奇点的解析类型（定理 6.15），因此翻模是奇点局部模型之间的无害操作。第二，终端化与 Q-分解可以同时完成，所得模型同时终端且 Q-分解，这就是三维 MMP 的标准工作平台。第三，三维标准翻模的存在由 $e(X)$ 的严格下降保证，而终止性则借助化归为射线型与定理 6.42 的有限性，两者合起来给出标准翻模序列的有限性。这三条结论说明：三维翻模的存在性与终止性并不需要预先假设 MMP 已经成立，这是本章相对于第 3 章的方法论进步。

辨析与提醒。第一，翻模与翻转不可混同：翻转沿收缩的标准类为负，翻模则满足 $K_X\equiv_f0$，两者都是小收缩，方向相反。第二，crepant 是相对于对 $(X,\Delta)$ 而言的，差异系数为 0 才叫 crepant，不能与终端、标准这两个奇点类别混为一谈。第三，本书完整证明的是终端翻模的终止（推论 6.46）；标准翻模只证明到简版（定理 6.43），一般情形引用 [Kol89]。阅读时须分清哪些结论是完整证明的、哪些是引用的。

评注。本章为第 7 章的标准翻模与标准模型铺路：§6.3 的终端 Q-分解模型是构造标准环的起点，§6.4 的存在性定理则保证标准翻模序列可以启动。历史上，三维翻转的存在性由森重文于 1988 年解决，而标准翻模的终止直到 Kollár 的工作才被完全证明；四维以上的对应问题至今仍是活跃的研究领域。读法建议：§6.1–6.2 的证明可先把握分次代数有限生成这条主线，§6.3 的符号幂论证较技术化，可在需要时回查。若时间有限，可先读 §6.1 与 §6.4，把存在性与终止性的主链条走一遍，再回头补 §6.3 的技术细节。}
